% A rising-side shape theorem. No assertion about signed products is made.
\section{Iterating positive differences of ordinary two-arm products}
Write $\mathcal D p=(\Delta p)^+$, with zero extension at the left.
Call a finite nonnegative sequence recursively rising-unimodal if every
$\mathcal D^j p$, $j\ge0$, is unimodal and log-concave on its rising
side, through its last mode.

\begin{lemma}
This class is preserved by terminal truncation and by convolution with
an interval polynomial $[t]$.
\end{lemma}
\begin{proof}
Truncation commutes with $\mathcal D$, apart from replacing the new
terminal negative difference by zero. Thus every iterated positive
difference is the corresponding truncation of the original one.
For windows, first apply the signed-window lemma to $\Delta p$:
its positive part is $\mathcal Dp$, and its positive support is an
initial interval. The window has a positive part which agrees with
$[t]\mathcal Dp$ until the end of the positive input, followed by a
nonincreasing part. Hence its next positive difference is a terminal
truncation of the corresponding positive difference of
$[t]\mathcal Dp$. Repeating this observation proves every order.
Equivalently one may induct on the order of differentiation; at each
step only a terminal truncation of a uniform window occurs. The
rising-side log-concavity at the joining edge follows from the
signed-window lemma. The assertion for order zero is the ordinary
rising-side log-concavity window lemma.
\end{proof}

\begin{theorem}
Let $A_T,A_W$ be formal arms whose ones occur initially and whose
remaining parts are at least two, in any order. Then the reverse of
$A_TA_W$ is recursively rising-unimodal. The analogous assertion
holds for a single reversed formal arm.
\end{theorem}
\begin{proof}
An interval is an immediate base case. A single reversed arm is
obtained by $q_{Tt}=[t]q_T+z^{\operatorname{wt}T+t}$.
The final two coefficients are both one, so this last monomial
does not change any positive difference. The window lemma proves
the assertion by induction.

For the product use its established unimodality and the established
strict first-mode bound $m_0<\max(a,b)$ when both factors are
nonrectangular. Orient $a\ge b$ and delete the last part $t$ of
the first arm. Then
\[
 p=[t]q_{T^-}q_W+z^a q_W.
\]
Through its first mode $m_0<a$, $p$ agrees with the window. Its
positive differences are therefore a terminal truncation of those
of that window; all further positive differences are truncated in
the same way. Apply induction and the window lemma. If a factor is
rectangular, the product is itself a uniform window and is covered
by the lemma. Rising-side log-concavity at order zero is the
already established ordinary-product assertion.
\end{proof}

In particular, the positive first differences are log-concave on
their own rising side. This does not assert their full log-concavity:
the previously recorded descending-side counterexamples still apply.
Nor does recursive rising-unimodality settle the signed two-arm
window required for the unrestricted one-descent problem.
